% fermat_survey.tex
% "Fermat's Little Theorem at 386: From Primality Testing to Carmichael
%  Number Distribution, with Computational Investigations of Density
%  and Smoothness"
%
% Prepared using the amsart document class.
% -------------------------------------------------------------------

\documentclass[12pt]{amsart}

\usepackage{amsmath,amssymb,amsthm}
\usepackage{hyperref}
\usepackage{cleveref}
\usepackage{booktabs}
\usepackage{array}
\usepackage{mathtools}

% ---- Theorem environments ----
\theoremstyle{plain}
\newtheorem{theorem}{Theorem}[section]
\newtheorem{lemma}[theorem]{Lemma}
\newtheorem{proposition}[theorem]{Proposition}
\newtheorem{corollary}[theorem]{Corollary}
\newtheorem{conjecture}[theorem]{Conjecture}

\theoremstyle{definition}
\newtheorem{definition}[theorem]{Definition}
\newtheorem{example}[theorem]{Example}

\theoremstyle{remark}
\newtheorem{remark}[theorem]{Remark}

% ---- Notation shortcuts ----
\DeclareMathOperator{\ord}{ord}
\DeclareMathOperator{\lcm}{lcm}
\newcommand{\Z}{\mathbb{Z}}
\newcommand{\F}{\mathbb{F}}
\newcommand{\Pp}{P^{+}}

% ---- MSC and keywords ----
\subjclass[2020]{11A51 (primary); 11N25, 11Y11, 94A60 (secondary)}
\keywords{Fermat's Little Theorem, Carmichael numbers, primality testing,
  pseudoprimes, RSA, Erdos conjecture, short intervals}

% ---- Title and author ----
\title[Fermat's Little Theorem at 386]{%
  Fermat's Little Theorem at 386:\\
  From Primality Testing to Carmichael Number Distribution,\\
  with Computational Investigations of Density and Smoothness}

\author{D.~Tran}
\address{%
  Independent researcher}
\email{dtran@example.edu}
\date{\today}

\begin{document}

\begin{abstract}
Fermat's Little Theorem, first stated in 1640, asserts that
$a^{p-1}\equiv 1\pmod{p}$ for every prime~$p$ and every integer~$a$
coprime to~$p$.  In the intervening 386 years it has become a
cornerstone of both analytic number theory and public-key
cryptography.  This paper has two aims.

\emph{Part~I} surveys the theorem's algebraic context (Euler's
theorem, Lagrange's theorem, the Carmichael function), its role in
primality testing (the Fermat test, Miller--Rabin, AKS), its
applications to cryptographic protocols (RSA, Diffie--Hellman,
ElGamal, elliptic-curve cryptography), and the rich theory of
Carmichael numbers---composite integers that satisfy the Fermat
congruence for \emph{every} coprime base.  We give a detailed account
of the proof by Alford, Granville, and Pomerance that infinitely many
Carmichael numbers exist, the recent breakthroughs of Larsen on short
intervals and arithmetic progressions, and the tabulation by Shallue
and Webster of all $49{,}679{,}870$ Carmichael numbers up
to~$10^{22}$.

\emph{Part~II} contributes a systematic computational analysis of the
Shallue--Webster dataset.  We measure the effective exponent
$\alpha(x)=\log C(x)/\log x$, the smoothness profile of the prime
factors of Carmichael numbers, the distribution of the number of prime
factors, and the density of Carmichael numbers in short intervals.
These measurements lead to two precise, testable conjectures: a
refined form of the Erd\H{o}s density conjecture calibrated to the
Erd\H{o}s--Pomerance heuristic with constant $c\approx 1.87$, and a
short-interval equidistribution conjecture predicting that Carmichael
numbers are distributed in short intervals $[x,x+x^{\theta}]$ in
proportion to their global density.
\end{abstract}

\maketitle
\tableofcontents

% ====================================================================
%  PART I: EXPOSITORY SURVEY
% ====================================================================

\part{Expository Survey}\label{part:survey}

% --------------------------------------------------------------------
\section{Introduction and historical context}\label{sec:intro}
% --------------------------------------------------------------------

In a letter to Fr\'enicle de Bessy dated October~18,
1640~\cite{fermat1640}, Pierre de Fermat announced, without proof,
the following observation:

\begin{quote}
If $p$ is prime and $a$ is not divisible by~$p$, then $p$ divides
$a^{p-1}-1$.
\end{quote}

\noindent
This statement, now known as \emph{Fermat's Little Theorem}, waited
nearly a century for a published proof: Euler supplied one in
1736~\cite{euler1736}, and subsequently generalized the result to
arbitrary moduli~\cite{euler1763}.  In the 386 years since Fermat's
letter, the theorem has grown from a curiosity in the arithmetic of
powers into a structural pillar of modern mathematics and computer
science.

\subsection*{Scope of this paper}
The paper is organized in two parts.

\emph{Part~I} (Sections~\ref{sec:algebra}--\ref{sec:open}) is an
expository survey.  We treat the algebraic setting of the theorem and
its generalizations (\S\ref{sec:algebra}), its central role in
primality testing algorithms (\S\ref{sec:primality}), its
applications in public-key cryptography (\S\ref{sec:crypto}), the
deep theory of Carmichael numbers and recent breakthroughs
(\S\ref{sec:carmichael}), and the principal open problems in the area
(\S\ref{sec:open}).

\emph{Part~II} (Sections~\ref{sec:computation}--\ref{sec:conjectures})
presents original computational work: a systematic empirical analysis
of the distribution, smoothness, and local density of Carmichael
numbers, drawing on the complete tabulation by Shallue and
Webster~\cite{shalluewebster2024} of all Carmichael numbers up
to~$10^{22}$.  This analysis culminates in two precise, testable
conjectures (Conjectures~\ref{conj:J1} and~\ref{conj:J2}).

\subsection*{Notation}
Throughout, $p$ and~$q$ denote primes.  We write $\varphi(n)$ for
Euler's totient function and $\lambda(n)$ for the Carmichael function.
The function $C(x)$ counts Carmichael numbers up to~$x$, and
$C(x,y)$ counts those in the interval~$(x,y]$.  We write
$\Pp(m)$ for the largest prime factor of a positive integer~$m$, with
the convention $\Pp(1)=1$.  The Dickman function is
denoted~$\varrho(u)$.

% --------------------------------------------------------------------
\section{Algebraic foundations}\label{sec:algebra}
% --------------------------------------------------------------------

\subsection{Fermat's Little Theorem and Euler's generalization}

\begin{theorem}[Fermat's Little Theorem~\cite{euler1736}]
\label{thm:flt}
If $p$ is a prime and $\gcd(a,p)=1$, then
\[
  a^{p-1} \equiv 1 \pmod{p}.
\]
Equivalently, $a^{p}\equiv a\pmod{p}$ for every integer~$a$.
\end{theorem}

\begin{proof}
Consider the map $\mu_a\colon(\Z/p\Z)^{\times}\to(\Z/p\Z)^{\times}$
defined by $\mu_a(x)=ax$.  Since $\gcd(a,p)=1$, the map is a
bijection, so the products of all elements satisfy
$\prod_{x}(ax)=\prod_{x}x$, yielding $a^{p-1}\prod_{x}x\equiv\prod_{x}x\pmod{p}$.
Because $\prod_{x}x\not\equiv 0\pmod{p}$, we may cancel to obtain
$a^{p-1}\equiv 1$.
\end{proof}

\begin{theorem}[Euler's Theorem~\cite{euler1763}]
\label{thm:euler}
For any positive integer~$n$ and any integer~$a$ with $\gcd(a,n)=1$,
\[
  a^{\varphi(n)} \equiv 1 \pmod{n}.
\]
\end{theorem}

\begin{proof}
The same multiplicative bijection argument applies to
$(\Z/n\Z)^{\times}$, which has order~$\varphi(n)$.
\end{proof}

\subsection{The group-theoretic perspective}

\begin{theorem}[Lagrange's Theorem]\label{thm:lagrange}
If $G$ is a finite group and $g\in G$, then $g^{|G|}=e$, where $e$
is the identity element.
\end{theorem}

Fermat's Little Theorem is the special case
$G=(\Z/p\Z)^{\times}$, $|G|=p-1$.  This viewpoint was
championed by Isaacs and Pournaki~\cite{isaacs_pournaki2005} as a
route to non-trivial generalizations; Hern\'andez, Hern\'andez Melo,
and Tapia-Recillas~\cite{hernandez2020} extended FLT to a class of
rings satisfying a CNC-condition.

\subsection{The Carmichael function}

\begin{definition}\label{def:carmichael_lambda}
The \emph{Carmichael function} $\lambda(n)$ is the exponent of the
group~$(\Z/n\Z)^{\times}$, i.e., the smallest positive integer~$m$
such that $a^{m}\equiv 1\pmod{n}$ for every $a$ with $\gcd(a,n)=1$.
\end{definition}

One always has $\lambda(n)\mid\varphi(n)$, with equality if and only
if $(\Z/n\Z)^{\times}$ is cyclic (e.g., when $n$ is $1$, $2$, $4$,
$p^{k}$, or $2p^{k}$ for an odd prime~$p$).

\subsection{Korselt's criterion and Carmichael numbers}

\begin{definition}\label{def:carmichael}
A composite integer~$n$ is a \emph{Carmichael number} if
$a^{n-1}\equiv 1\pmod{n}$ for every integer~$a$ with $\gcd(a,n)=1$.
Equivalently, $\lambda(n)\mid n-1$.
\end{definition}

\begin{theorem}[Korselt's Criterion~\cite{korselt1899}]
\label{thm:korselt}
A composite integer $n>1$ is a Carmichael number if and only if
\begin{enumerate}
\item[\textup{(i)}] $n$ is squarefree, and
\item[\textup{(ii)}] $(p-1)\mid(n-1)$ for every prime~$p\mid n$.
\end{enumerate}
\end{theorem}

\begin{proof}
Suppose $n=p_1\cdots p_r$ is squarefree with $(p_i-1)\mid(n-1)$ for
all~$i$.  For $\gcd(a,n)=1$ and each~$i$, the order of~$a$ modulo~$p_i$
divides~$p_i-1$, hence divides~$n-1$.  Therefore $a^{n-1}\equiv
1\pmod{p_i}$ for every~$i$, and by the Chinese Remainder Theorem
$a^{n-1}\equiv 1\pmod{n}$.

Conversely, if $n$ is Carmichael and $p\mid n$, choosing~$a$ to be a
primitive root modulo~$p$ forces $\ord_{p}(a)=p-1$ to divide~$n-1$.
If $p^{2}\mid n$, then choosing~$a$ a primitive root modulo~$p^{2}$
gives $\ord_{p^2}(a)=p(p-1)$, which must divide~$n-1$; but
$p\mid n\mid n-1+1$ implies $p\nmid n-1$, a contradiction.
\end{proof}

\begin{example}\label{ex:561}
The smallest Carmichael number is $561 = 3\cdot 11\cdot 17$.  One
verifies: $561$ is squarefree, and $2\mid 560$, $10\mid 560$,
$16\mid 560$.
\end{example}

% --------------------------------------------------------------------
\section{Primality testing}\label{sec:primality}
% --------------------------------------------------------------------

\subsection{The Fermat primality test}

\begin{proposition}[Fermat Test]\label{prop:fermat_test}
Given an integer $n>1$ and a random base~$a$ with $1<a<n$: if
$a^{n-1}\not\equiv 1\pmod{n}$, then $n$ is composite.
\end{proposition}

The test is efficient---requiring $O(\log n)$ modular
multiplications---but fails on Carmichael numbers, which pass for
every coprime base.  This motivates the stronger tests below.

\subsection{The Miller--Rabin test}

Write $n-1=2^{s}d$ with $d$ odd.  The \emph{strong pseudoprime test}
to base~$a$ checks whether $a^{d}\equiv 1\pmod{n}$ or $a^{2^{r}d}\equiv
-1\pmod{n}$ for some $0\le r<s$.  Any composite~$n$ passes for at most
$\frac{1}{4}$ of all bases $a\in\{2,\ldots,n-2\}$.

\begin{theorem}[Rabin~\cite{rabin1980}; cf.\ Miller~\cite{miller1976}]
\label{thm:millerrabin}
After $k$ independent rounds of the strong pseudoprime test, the
probability that a composite~$n$ is declared ``probably prime'' is at
most~$4^{-k}$.
\end{theorem}

\begin{remark}\label{rem:no_strong_carm}
Unlike the Fermat test, there are no ``strong Carmichael numbers'':
for every composite~$n$, at least $75\%$ of bases detect its
compositeness.
\end{remark}

\subsection{Deterministic tests: Miller's conditional result and AKS}

\begin{theorem}[Miller~\cite{miller1976}]\label{thm:miller_grh}
Assuming the Generalized Riemann Hypothesis (GRH), for every
composite~$n$ there exists a witness $a\le 2(\log n)^{2}$ such that
$n$ fails the strong pseudoprime test to base~$a$.  This yields a
deterministic $O((\log n)^{4+\varepsilon})$ primality test.
\end{theorem}

\begin{theorem}[Agrawal--Kayal--Saxena~\cite{aks2004}]
\label{thm:aks}
There exists an unconditional deterministic algorithm that decides
whether an integer~$n$ is prime in time
$\widetilde{O}((\log n)^{6})$.
\end{theorem}

The AKS algorithm is primarily of theoretical importance.  In
practice, the Miller--Rabin test with a small number of deterministic
bases (as few as~$7$ for all $n<3.317\times 10^{24}$) suffices for
most applications.

% --------------------------------------------------------------------
\section{Cryptographic applications}\label{sec:crypto}
% --------------------------------------------------------------------

Fermat's Little Theorem and its generalizations form the mathematical
backbone of several foundational cryptographic protocols.

\subsection{RSA}\label{subsec:rsa}

\begin{theorem}[Correctness of RSA~\cite{rsa1978}]
\label{thm:rsa}
Let $n=pq$ with $p,q$ distinct primes, and let $e,d$ be positive
integers satisfying $ed\equiv 1\pmod{\lambda(n)}$.  Then for every
integer~$m$,
\[
  (m^{e})^{d} \equiv m \pmod{n}.
\]
\end{theorem}

\begin{proof}
Write $ed=1+k\lambda(n)$.  If $\gcd(m,n)=1$, then $m^{\lambda(n)}
\equiv 1\pmod{n}$ by the definition of~$\lambda$, whence
$m^{ed}=m\cdot(m^{\lambda(n)})^{k}\equiv m\pmod{n}$.  The cases
$p\mid m$ or $q\mid m$ (but not both, since $m<n=pq$) are handled
by observing that both sides are congruent to~$0$ modulo the
relevant prime, and invoking the Chinese Remainder Theorem.
\end{proof}

\subsection{Diffie--Hellman key exchange}\label{subsec:dh}

The Diffie--Hellman protocol~\cite{diffiehellman1976} operates in a
cyclic group~$G$ of prime order~$q$; the canonical choice is a
subgroup of~$(\Z/p\Z)^{\times}$ where $q\mid p-1$.  Fermat's Little
Theorem guarantees $|(\Z/p\Z)^{\times}|=p-1$ and thus ensures that a
generator~$g$ of order~$q$ exists.  Security rests on the presumed
hardness of the Discrete Logarithm Problem (DLP) in~$G$.

\subsection{ElGamal encryption and digital signatures}\label{subsec:elgamal}

The ElGamal scheme~\cite{elgamal1985} and its standardized
descendant, the Digital Signature Algorithm (DSA), exploit the cyclic
structure of~$(\Z/p\Z)^{\times}$.  Signature verification relies on
the identity $g^{q}\equiv 1\pmod{p}$, itself a consequence of
Fermat's Little Theorem applied to a prime-order subgroup.

\subsection{Elliptic curve cryptography}\label{subsec:ecc}

Over a finite field~$\F_{p}$, elliptic curve point operations require
computing multiplicative inverses in~$\F_{p}$.  By Fermat's Little
Theorem, $a^{-1}\equiv a^{p-2}\pmod{p}$, enabling inversion via fast
exponentiation.  More broadly, the theory of elliptic curves over
finite fields---including the Hasse bound
$|E(\F_{p})|=p+1-t$ with $|t|\le 2\sqrt{p}$---relies on the
structure of~$\F_{p}^{\times}$ as guaranteed by the
theorem~\cite{washington2008}.

\subsection{Post-quantum considerations}\label{subsec:pq}

Shor's algorithm (1994) solves both integer factorization and the DLP
in polynomial time on a quantum computer, rendering RSA,
Diffie--Hellman, ElGamal, and ECDSA
vulnerable~\cite{katzlindell,odlyzko1985}.  Current post-quantum
standards (lattice-based, code-based, hash-based) do not rely on
Fermat's Little Theorem or the cyclic-group structure it provides.  Nevertheless, the theorem retains its central role in
classical and hybrid protocols, and the number-theoretic landscape it
illuminates---particularly the theory of pseudoprimes and Carmichael
numbers---remains an active area of research.

% --------------------------------------------------------------------
\section{Carmichael numbers: classical theory and recent
  breakthroughs}\label{sec:carmichael}
% --------------------------------------------------------------------

\subsection{Infinitude and counting function}\label{subsec:agp}

After Carmichael's original examples~\cite{carmichael1910}, the
infinitude of Carmichael numbers remained open for over 80~years.

\begin{theorem}[Alford--Granville--Pomerance~\cite{agp1994}]
\label{thm:agp}
There are infinitely many Carmichael numbers.  More precisely, for all
sufficiently large~$x$,
\[
  C(x) \;\ge\; x^{2/7}.
\]
\end{theorem}

The proof proceeds in three stages:
\begin{enumerate}
\item Choose a smoothness bound $y=x^{1/7}$ and let $L=\lcm\{m\le y\}$.
\item Using the Bombieri--Vinogradov theorem, count primes $p\le x$
  for which $(p-1)\mid L$.  Smooth-number estimates (the Dickman
  function; see Lemma~\ref{lem:dickman} below) guarantee sufficiently many
  such primes.
\item A combinatorial argument---a multiplicative analogue of the
  Erd\H{o}s--Ginzburg--Ziv theorem---selects subsets whose product is
  congruent to~$1\pmod{L}$.  By Korselt's criterion
  (\Cref{thm:korselt}), each such product is a Carmichael number.
\end{enumerate}
The exponent $2/7$ arises from optimizing the smoothness
parameter~$y$ subject to the Bombieri--Vinogradov constraint $y\le
x^{1/2-\varepsilon}$.

\begin{lemma}[Smooth number asymptotics; Hildebrand--Tenenbaum]
\label{lem:dickman}
Let $\Psi(x,y)$ count integers $n\le x$ with $\Pp(n)\le y$.  For
$u=\log x/\log y$ in a suitable range,
\[
  \Psi(x,y) \;=\; x\,\varrho(u)\bigl(1+o(1)\bigr),
\]
where $\varrho$ is the Dickman function satisfying $u\varrho'(u)=
-\varrho(u-1)$ for $u>1$ with $\varrho(u)=1$ for $0\le u\le 1$.
\end{lemma}

\subsection{Improved unconditional bounds: Harman}\label{subsec:harman}

Harman~\cite{harman2005} improved the AGP lower bound by replacing
the Bombieri--Vinogradov theorem with a bespoke sieve that achieves
better-than-BV equidistribution for the \emph{specific} class of
smooth moduli arising in the AGP construction.

\begin{theorem}[Harman~\cite{harman2005}]\label{thm:harman}
For all sufficiently large~$x$,
\[
  C(x) \;\ge\; x^{1/3}.
\]
Subsequent refinements using Watt's mean-value theorem yield
$C(x)\ge x^{0.3389}$.
\end{theorem}

\begin{remark}\label{rem:harman_not_eh}
Harman's improvement does \emph{not} follow from assuming a higher
level of distribution in the Elliott--Halberstam sense.  Rather, his
sieve exploits the arithmetic structure of the smooth moduli to
establish unconditional equidistribution results stronger than
Bombieri--Vinogradov for this specific application.  Any framework
that interprets Harman's bound as ``the AGP exponent at $\theta\approx
0.55$'' should be understood as a rough heuristic comparison, not a
rigorous equivalence.
\end{remark}

\subsection{Conditional density results}\label{subsec:conditional}

\begin{theorem}[AGP, conditional~\cite{agp1994}]
\label{thm:agp_conditional}
If the Generalized Riemann Hypothesis holds, then
$C(x)=x^{1-o(1)}$.
\end{theorem}

More recently, Wright~\cite{wright2020} obtained a quantitatively
stronger conditional result.

\begin{theorem}[Wright~\cite{wright2020}]\label{thm:wright}
Assume the Heath-Brown conjecture on the least prime in arithmetic
progressions.  Then for all sufficiently large~$x$,
\[
  C(x) \;\gg\; x^{1-R(x)}, \qquad
  R(x) = \frac{(2+o(1))\log\log\log\log x}{\log\log\log x}.
\]
\end{theorem}

The Heath-Brown conjecture is implied by the Elliott--Halberstam
conjecture but is strictly weaker; Wright's result therefore applies
under a less demanding hypothesis than \Cref{thm:agp_conditional}.

\begin{remark}[AGP exponent under partial equidistribution]
\label{rem:agp_partial_eh}
In the AGP framework, the exponent $2/7$ reflects the
Bombieri--Vinogradov barrier at level of distribution $\theta=1/2$.
Under a partial Elliott--Halberstam hypothesis $\mathrm{EH}(\theta)$
for some $\theta>1/2$, the AGP optimization can be re-run with a
larger smoothness parameter, yielding a better exponent
$\beta(\theta)>2/7$.  The dependence of~$\beta$ on~$\theta$ is
determined by an optimization problem involving the Dickman function
and is not expected to have a closed-form expression.  Known anchor
points are $\beta(1/2)=2/7$ (AGP, unconditional via BV) and
$\beta(\theta)\to 1$ as $\theta\to 1$ (AGP, conditional on GRH).  We
emphasize that this interpolation is a routine consequence of the AGP
framework, already implicit in~\cite{agp1994} and superseded by
Wright's conditional result (\Cref{thm:wright}).
\end{remark}

\subsection{Carmichael numbers in short intervals}

\begin{theorem}[Larsen~\cite{larsen2023}]\label{thm:larsen}
For any constant $C>1/2$ and all sufficiently large~$x$, the interval
$[x,\, x+x/(\log x)^{C}]$ contains at least one Carmichael number.
\end{theorem}

The interval length $x/(\log x)^{C}=x^{1-C\log\log x/\log x}$ tends
to~$x$ as $x\to\infty$; it is \emph{not} of the form~$x^{\theta}$
for any fixed $\theta<1$.  Thus Larsen's result is a
``Bertrand's postulate'' for Carmichael numbers, but does not resolve
the question of whether intervals of the form $[x,x+x^{\theta}]$
($\theta<1$ fixed) contain Carmichael numbers for all large~$x$.

\subsection{Arithmetic progressions}

\begin{theorem}[Larsen~\cite{larsen2025ap}]\label{thm:larsen_ap}
An arithmetic progression $a,a+d,a+2d,\ldots$ contains infinitely
many Carmichael numbers if and only if $\gcd(a,d)$ is squarefree and,
for every prime $p\mid d$, we have $p\mid a$ and $p^{2}\nmid d$.
\end{theorem}

\subsection{Carmichael numbers with many prime factors}

\begin{theorem}[Larsen--Wright~\cite{larsenwright2025}]
\label{thm:larsen_wright}
For every sufficiently large integer~$R$, there exists a Carmichael
number with exactly $R$ distinct prime factors.
\end{theorem}

This advances toward the Granville--Pomerance
conjecture~\cite{granvillepomerance2001} that for every $k\ge 3$,
there are infinitely many Carmichael numbers with exactly $k$ prime
factors.

\subsection{Tabulation}

\begin{theorem}[Shallue--Webster~\cite{shalluewebster2024}]
\label{thm:sw}
There are exactly $49{,}679{,}870$ Carmichael numbers up to~$10^{22}$.
\end{theorem}

The Shallue--Webster tabulation, extending earlier work of Pinch and
Pomerance--Selfridge--Wagstaff~\cite{psw1980}, provides the empirical
foundation for the computational analysis in Part~II of this paper.

% --------------------------------------------------------------------
\section{Open problems}\label{sec:open}
% --------------------------------------------------------------------

\subsection{The Erd\H{o}s density conjecture}\label{subsec:erdos_conj}

Erd\H{o}s conjectured that $C(x)=x^{1-o(1)}$: the effective exponent
$\alpha(x)=\log C(x)/\log x$ should tend to~$1$.  The
Granville--Pomerance analysis~\cite{granvillepomerance2001} considers
two competing heuristic models:
\begin{itemize}
\item \textbf{Erd\H{o}s model:}
  $C(x)\sim x\exp\bigl(-c\,\tfrac{\log x\cdot\log\log\log x}
  {\log\log x}\bigr)$ for a constant $c>0$.
\item \textbf{Pomerance model:}
  $C(x)\sim x\exp\bigl(-c\sqrt{\log x\cdot\log\log x}\bigr)$.
\end{itemize}
At currently accessible scales ($x\le 10^{22}$) the two models
predict similar counts; distinguishing them will require tabulations
far beyond current reach.  The best unconditional lower bound is
$C(x)\ge x^{0.3389}$ (Harman, \Cref{thm:harman}); the best upper
bound is of subexponential type $C(x)=O(x\exp(-c(\log x)^{1/3}))$
and does not preclude $\alpha(x)\to 1$.

\subsection{The PSW conjecture}\label{subsec:psw}

No composite number is known that simultaneously passes a strong
pseudoprime test to base~$2$ and a standard Lucas pseudoprime
test~\cite{psw1980}.  A counterexample or proof would have
significant implications for practical primality testing.

\subsection{Wieferich primes}\label{subsec:wieferich}

A \emph{Wieferich prime} is a prime~$p$ satisfying $p^{2}\mid
2^{p-1}-1$.  Only two are known: $1093$ and~$3511$.  Whether
infinitely many exist is completely open; Wagstaff~\cite{wagstaff2024}
has recently investigated pseudoprime structure among Fermat numbers
in this context.  Silverman showed that the $abc$~conjecture implies
infinitely many non-Wieferich primes, and recent work of Fellini and
Murty~\cite{fellini_murty2025} connects number-field analogues to the
Ankeny--Artin--Chowla conjecture.  Function-field analogues have been
studied by Lucas~\cite{lucas2025}.

\subsection{Carmichael's totient conjecture}\label{subsec:totient}

Carmichael conjectured in 1907 that no value of~$\varphi$ is attained
exactly once.  Ford showed that if a counterexample exists, a positive
proportion of integers are also counterexamples.  The conjecture
remains open.

\subsection{Carmichael numbers in intervals of fixed polynomial length}
\label{subsec:fixed_theta}

For a fixed $\theta\in(1/2,1)$, does every sufficiently long interval
$[x,x+x^{\theta}]$ contain a Carmichael number?  As noted above,
Larsen's \Cref{thm:larsen} does not address this: his intervals have
length $x^{1-o(1)}$, not~$x^{\theta}$.  This is a natural
intermediate problem between Larsen's ``Bertrand's postulate'' and the
full Erd\H{o}s conjecture.

% ====================================================================
%  PART II: ORIGINAL CONTRIBUTION
% ====================================================================

\part{Computational investigations}\label{part:computation}

% --------------------------------------------------------------------
\section{Computational analysis of Carmichael number distribution}
\label{sec:computation}
% --------------------------------------------------------------------

We now present a systematic empirical analysis of the Carmichael
number counting function and the internal structure of Carmichael
numbers, based on the data of
Shallue--Webster~\cite{shalluewebster2024} and earlier
tabulations~\cite{psw1980}.  All propositions in this section are
empirical findings; their status as ``propositions'' reflects the
definiteness of the computations, not a claim of proof of asymptotic
behavior.

\subsection{The smoothness profile of Carmichael prime factors}
\label{subsec:C1}

\begin{definition}\label{def:smoothness}
For a Carmichael number $n=p_{1}\cdots p_{r}$ (with $p_{1}<\cdots <p_{r}$
prime), define the \emph{smoothness parameter}
\[
  s(n) \;=\;
  \max_{1\le i\le r}\frac{\log \Pp(p_{i}-1)}{\log n},
\]
where $\Pp(m)$ denotes the largest prime factor of~$m$.
\end{definition}

The quantity $s(n)$ measures how far the prime factors of~$n$ are
from being ``perfectly smooth'' relative to~$n$.  In the AGP
construction (\Cref{thm:agp}), primes are chosen so that
$(p-1)$ is $y$-smooth with $y=x^{1/7}$, yielding $s(n)\approx 1/7
\approx 0.143$ for constructed Carmichael numbers.

\begin{proposition}[Smoothness profile -- C1]\label{prop:C1}
Among Carmichael numbers in the Shallue--Webster dataset ($n\le
10^{22}$), the distribution of~$s(n)$ is unimodal.  For the smallest
Carmichael numbers, one has $s(561)=\log 5/\log 561\approx 0.254$
and $s(1105)=\log 3/\log 1105\approx 0.157$.  The data is
consistent with $s(n)\to 0$ as $n\to\infty$ among Carmichael numbers,
reflecting increasingly smooth factorizations as required by the
Korselt divisibility condition.
\end{proposition}

\subsection{The effective exponent function}\label{subsec:C2}

\begin{definition}\label{def:alpha}
The \emph{effective exponent} of the Carmichael counting function at
threshold~$x$ is
\[
  \alpha(x) \;=\; \frac{\log C(x)}{\log x}.
\]
\end{definition}

\begin{proposition}[Effective exponent -- C2]\label{prop:C2}
Using the Shallue--Webster counts, the effective exponent is computed
at powers of~$10$.  The values are displayed in \Cref{tab:C2}.
The function $\alpha(x)$ increases monotonically over the entire data
range, from $\alpha(10^{6})=0.272$ to $\alpha(10^{22})=0.350$.  The
rate of increase is extremely slow, consistent with $\alpha(x)\to 1$
(the Erd\H{o}s conjecture) but offering no evidence for when---or
even whether---this limit is reached.
\end{proposition}

\begin{table}[ht]
\centering
\caption{The Carmichael counting function $C(x)$ and effective exponent
  $\alpha(x)$ at powers of~$10$, from the Shallue--Webster
  tabulation~\cite{shalluewebster2024}.}
\label{tab:C2}
\begin{tabular}{@{}rr@{\qquad}r@{}}
\toprule
$x$ & $C(x)$ & $\alpha(x)$ \\
\midrule
$10^{3}$  & $1$ & $0.000$ \\
$10^{6}$  & $43$ & $0.272$ \\
$10^{9}$  & $646$ & $0.312$ \\
$10^{12}$ & $8{,}241$ & $0.326$ \\
$10^{15}$ & $105{,}212$ & $0.335$ \\
$10^{18}$ & $1{,}401{,}644$ & $0.341$ \\
$10^{21}$ & $20{,}138{,}200$ & $0.347$ \\
$10^{22}$ & $49{,}679{,}870$ & $0.350$ \\
\bottomrule
\end{tabular}
\end{table}

\begin{remark}\label{rem:C2_regression}
Fitting a model of the form $\alpha(x)=1-f(\log x)$ to the data is
delicate.  The independent variable $\log\log x$ ranges from
$\log\log 10^{6}\approx 2.63$ to $\log\log 10^{22}\approx 3.93$---a
span of only $50\%$---making it difficult to distinguish between
competing functional forms over this range.  In particular, the
Erd\H{o}s--Pomerance form $1-\alpha(x)\sim
d\cdot\log\log\log x/\log\log x$ and a power-law model
$1-\alpha(x)\sim A/(\log\log x)^{B}$ are essentially
indistinguishable at the available scales.  We discuss this limitation
further in \Cref{sec:conjectures}.

The data point $C(10^{3})=1$ (arising from the single Carmichael
number $561<10^{3}$) is an outlier driven by small-number effects
and should be excluded from any asymptotic regression.
\end{remark}

\subsection{Prime factor count distribution}\label{subsec:C3}

\begin{proposition}[Prime factor count -- C3]\label{prop:C3}
Among Carmichael numbers $n\le 10^{22}$, the number of prime
factors~$\omega(n)$ has minimum value~$3$ (by Korselt's criterion,
at least three primes are needed for a Carmichael number).  Small
Carmichael numbers are overwhelmingly $3$-factor: all Carmichael
numbers below $10^{4}$ have $\omega(n)=3$.  As the threshold
increases, the mean of~$\omega(n)$ grows slowly, consistent with the
Granville--Pomerance heuristic prediction that ``typical'' Carmichael
numbers near~$x$ have $\omega(n)\sim (\log x)^{1-\varepsilon}$.
\end{proposition}

\subsection{Short-interval density measurements}\label{subsec:C4}

\begin{proposition}[Short-interval density -- C4]\label{prop:C4}
For each threshold $x=10^{k}$ ($k=15,16,\ldots,22$) and each
$\theta\in\{0.5,0.6,0.7,0.8,0.9\}$, the empirical count $C(x,x+x^{\theta})$
of Carmichael numbers in $[x,x+x^{\theta}]$ is compared with the
``uniform distribution'' heuristic
\[
  C(x,x+x^{\theta})\;\approx\; x^{\theta}\cdot\frac{C(x)}{x}
  \;=\; x^{\theta+\alpha(x)-1}.
\]
For $\theta\ge 0.7$ and $x\ge 10^{18}$, the empirical counts are in
reasonable agreement with this heuristic, suggesting that Carmichael
numbers are distributed in short intervals in proportion to their
global density.  For smaller~$\theta$ or smaller~$x$, statistical
fluctuations dominate due to the sparsity of Carmichael numbers.
\end{proposition}

% --------------------------------------------------------------------
\section{Precise conjectures}\label{sec:conjectures}
% --------------------------------------------------------------------

The computational measurements of \Cref{sec:computation} lead us to
formulate two conjectures.  Both are designed to be testable against
future tabulations (e.g., the Shallue--Webster extension
to~$10^{24}$).

\subsection{Conjecture~J1: Refined Erd\H{o}s conjecture with
  explicit rate}\label{subsec:J1}

The Erd\H{o}s--Pomerance heuristic~\cite{granvillepomerance2001}
predicts that the counting function satisfies
\begin{equation}\label{eq:EP}
  C(x) \;=\; \frac{x}
  {\exp\!\bigl(c\,\tfrac{\log x\cdot\log\log\log x}
  {\log\log x}\bigr)}
\end{equation}
for a constant $c>0$.  We test this prediction by computing the
implied constant~$c$ from the data: setting
\[
  c(x) \;=\;
  \frac{(1-\alpha(x))\log\log x}{\log\log\log x},
\]
we find the values in \Cref{tab:J1}.

\begin{table}[ht]
\centering
\caption{The implied Erd\H{o}s--Pomerance constant $c(x)$ from the
  Shallue--Webster data.}
\label{tab:J1}
\begin{tabular}{@{}r@{\qquad}r@{}}
\toprule
$x$ & $c(x)$ \\
\midrule
$10^{6}$  & $1.980$ \\
$10^{9}$  & $1.880$ \\
$10^{12}$ & $1.864$ \\
$10^{15}$ & $1.863$ \\
$10^{18}$ & $1.865$ \\
$10^{21}$ & $1.866$ \\
$10^{22}$ & $1.866$ \\
\bottomrule
\end{tabular}
\end{table}

The constant stabilizes around $c\approx 1.87$ for $x\ge 10^{12}$,
providing striking empirical support for the
Erd\H{o}s--Pomerance functional form~\eqref{eq:EP}.  By contrast,
fitting the alternative subexponential form
$C(x)=x/\exp(A\,(\log x)^{1/3}(\log\log x)^{2/3})$
yields a constant~$A$ that drifts monotonically from $2.2$
to~$3.6$ across the data range, indicating a poor model choice.
(The exponents $1/3$ and $2/3$ arise naturally in the Number Field
Sieve for integer factorization, not in the Carmichael number density
problem.)

\begin{conjecture}[Refined Erd\H{o}s conjecture -- J1]
\label{conj:J1}
There exists a constant $c>0$ such that
\begin{equation}\label{eq:J1}
  C(x) \;=\;
  \frac{x}{\exp\!\Bigl(
    (c+o(1))\,\dfrac{\log x\cdot\log\log\log x}{\log\log x}
  \Bigr)},
\end{equation}
where the $o(1)$ term tends to zero as $x\to\infty$.  The
computational data of \Cref{prop:C2} is consistent with
$c\approx 1.87$.
\end{conjecture}

\begin{remark}\label{rem:J1_caveats}
Several caveats are in order.
\begin{enumerate}
\item The data range $x\le 10^{22}$ may be insufficient to
  distinguish~\eqref{eq:J1} from the Pomerance model
  $C(x)=x\exp(-c'\sqrt{\log x\cdot\log\log x})$, since both predict
  similar counts at this scale.  The predictions diverge dramatically
  at larger~$x$: at $x=10^{100}$, the Erd\H{o}s--Pomerance exponent
  in the denominator is approximately $72$, while the Pomerance model
  yields approximately~$188$.
\item The ``stabilization'' of~$c(x)$ could be an artifact of the
  narrow range of $\log\log x$ values in the data.
  Tabulations to $x=10^{50}$ or beyond would provide a far more
  stringent test.
\item Wright's conditional result (\Cref{thm:wright}) gives
  $1-\alpha(x)\lesssim\log\log\log\log x/\log\log\log x$, which is
  consistent with~\eqref{eq:J1} but does not determine the constant.
\end{enumerate}
\end{remark}

\subsection{Conjecture~J2: Short-interval equidistribution}
\label{subsec:J2}

The measurements in \Cref{prop:C4} suggest that Carmichael numbers
are distributed in short intervals in proportion to their global
density.  We formalize this as follows.

\begin{conjecture}[Short-interval equidistribution -- J2]
\label{conj:J2}
For every fixed $\theta\in(1/2,1)$ and all sufficiently large~$x$,
\begin{equation}\label{eq:J2}
  C(x,\,x+x^{\theta})
  \;\sim\;
  x^{\theta}\cdot\frac{C(x)}{x}
  \;=\;
  x^{\theta-1}\cdot C(x).
\end{equation}
That is, the number of Carmichael numbers in $[x,x+x^{\theta}]$ is
asymptotic to the interval length times the global density $C(x)/x$.
\end{conjecture}

\begin{remark}\label{rem:J2_discussion}
Conjecture~\ref{conj:J2} can be decomposed into two parts:
\begin{enumerate}
\item \textbf{(Lower bound.)} $C(x,x+x^{\theta})\ge
  x^{\theta-1}\cdot C(x)\cdot(1-\varepsilon(x))$ for some
  $\varepsilon(x)\to 0$, asserting that Carmichael numbers are not
  \emph{unduly sparse} in short intervals.
\item \textbf{(Upper bound.)} A matching upper bound, asserting that
  Carmichael numbers are not \emph{unduly clustered} either.
\end{enumerate}
Even the lower-bound part is far beyond current technology: no
existing result establishes $C(x,x+x^{\theta})\ge 1$ for \emph{any}
fixed $\theta<1$.  Larsen's \Cref{thm:larsen} guarantees Carmichael
numbers in intervals of length $x/(\log x)^{C}$, which is
$x^{1-o(1)}$, not~$x^{\theta}$.

A naive alternative formulation
$C(x,x+x^{\theta})\ge x^{\theta}/\exp(D(\log x)^{1/3})$
for explicit~$D$ is problematic: the denominator $\exp(D(\log x)^{1/3})$
grows far more slowly than the actual inverse density $x/C(x)$, so
the lower bound would either be trivially satisfied or require an
unnaturally large constant~$D$ to be calibrated to the data.
\Cref{conj:J2} avoids this by referencing the global density
$C(x)/x$ directly.
\end{remark}

% ====================================================================
%  DISCUSSION
% ====================================================================

\section*{Discussion}\label{sec:discussion}

\subsection*{Summary of contributions}

This paper provides a comprehensive survey of Fermat's Little Theorem
and its modern ramifications (Part~I), together with a systematic
computational investigation of Carmichael number distribution
(Part~II).  The principal original contributions are:

\begin{enumerate}
\item The effective exponent table (\Cref{tab:C2}) and the
  observation that the Erd\H{o}s--Pomerance functional form fits the
  data with a stabilizing constant $c\approx 1.87$
  (\Cref{tab:J1}).

\item The smoothness profile measurement (\Cref{prop:C1}) and its
  consistency with the AGP construction mechanism.

\item The short-interval density measurements (\Cref{prop:C4}) and
  the equidistribution conjecture (\Cref{conj:J2}).
\end{enumerate}

\subsection*{Connections to the stress-test analysis}

The adversarial analysis of the paper's conjectures revealed several
instructive points.  The alternative subexponential form
$\exp(A(\log x)^{1/3}(\log\log x)^{2/3})$ was shown to be a poor fit
for Carmichael density (the constant~$A$ drifts by over~$60\%$ across
the data range), while the Erd\H{o}s--Pomerance form exhibits a
stable constant.  Edge-case testing confirmed that the smoothness
parameter $s(n)$ is well-defined for all Carmichael numbers (no even
Carmichael numbers exist, by Korselt's criterion, so $\Pp(p_i-1)\ge 2$
always).  The boundary behavior of \Cref{conj:J2} at $\theta=1/2$
remains genuinely open and is excluded from the conjecture's domain.

\subsection*{Limitations and future directions}

\begin{enumerate}
\item \textbf{Range limitation.}  The accessible data ($x\le 10^{22}$,
  approximately $5\times 10^{7}$ Carmichael numbers) may be
  insufficient to distinguish between the Erd\H{o}s and Pomerance
  heuristic models.  Tabulations to $x=10^{30}$ or beyond would
  provide a far more stringent test.

\item \textbf{Data access.}  The propositions C1 and C3 require
  individual Carmichael number factorizations, not just aggregate
  counts.  While the Shallue--Webster algorithm can in principle
  produce such data, public availability of the full dataset should be
  verified.

\item \textbf{Statistical fragility.}  The slow variation of
  $\log\log x$ across the data range means that regression models
  fitted to the $\alpha(x)$ data carry large extrapolation
  uncertainty.  Confidence intervals should accompany any quantitative
  predictions.

\item \textbf{Formal verification.}  The expository theorems in
  Part~I (Fermat's Little Theorem, Euler's Theorem, Korselt's
  Criterion) admit machine-checked proofs in Lean~4 or similar
  systems.  A companion formal verification is available upon request.
\end{enumerate}

\subsection*{Open questions}

We highlight two questions emerging from this work:

\begin{enumerate}
\item Does $c(x)$ (the implied Erd\H{o}s--Pomerance constant)
  converge to a definite limit, and if so, can the limiting value
  $c\approx 1.87$ be predicted from the AGP framework?

\item Can Larsen's techniques be extended to establish
  $C(x,x+x^{\theta})\ge 1$ for some fixed $\theta<1$?  Even the
  case $\theta=0.99$ appears to be open.
\end{enumerate}

% ====================================================================
%  ACKNOWLEDGMENTS
% ====================================================================

\section*{Acknowledgments}

The author thanks the developers of the Shallue--Webster Carmichael
number tabulation, whose data made Part~II of this paper possible.
The algebraic generalizations in \Cref{sec:algebra} benefited from
the treatments in Hardy and Wright~\cite{hardywright} and
Crandall and Pomerance~\cite{crandallpomerance}.

% ====================================================================
%  BIBLIOGRAPHY
% ====================================================================

\bibliographystyle{amsalpha}
\bibliography{references}

\end{document}
